\section{Choix du modèle de base}

Pour cette expérimentation, nous avons choisi \textbf{MobileNetV2} \cite{sandler2018mobilenetv2} comme architecture de base pour notre système de détection. MobileNetV2 est un réseau de neurones convolutif léger développé par Google, caractérisé par une architecture basée sur les \textit{inverted residual blocks} et les \textit{linear bottlenecks}. Ces innovations architecturales permettent d'obtenir des performances de classification compétitives (72\% top-1 accuracy sur ImageNet) tout en maintenant une empreinte computationnelle très réduite : seulement 3.5 millions de paramètres et 300 millions de FLOPs, contre 25.6 millions de paramètres et 4.1 milliards de FLOPs pour ResNet50. Le modèle accepte des images de taille $224 \times 224 \times 3$ en entrée et utilise des convolutions depthwise séparables pour minimiser le coût computationnel tout en préservant la capacité de représentation.

Le choix de MobileNetV2 repose sur trois critères principaux. Premièrement, son \textbf{efficacité computationnelle} (14 MB de taille, temps d'inférence de 25ms sur CPU) permet un déploiement en production sur des environnements aux ressources limitées, sans nécessiter de GPU dédié. Deuxièmement, sa \textbf{structure modulaire} avec bottlenecks linéaires se prête particulièrement bien à l'adversarial training : la faible complexité du modèle réduit le temps de génération des exemples adverses (FGSM, PGD) à chaque batch, accélérant significativement la convergence durant l'entraînement robuste. Troisièmement, l'architecture offre un excellent support pour le \textbf{transfer learning} : les poids pré-entraînés sur ImageNet fournissent des features génériques de haute qualité que nous adaptons à notre tâche de détection d'armes à feu par fine-tuning des dernières couches, tout en conservant la représentation des patterns visuels fondamentaux (textures, contours, formes) appris sur le dataset source.



Référence bibliographique à ajouter dans le fichier .bib :
@inproceedings{sandler2018mobilenetv2,
  title={MobileNetV2: Inverted Residuals and Linear Bottlenecks},
  author={Sandler, Mark and Howard, Andrew and Zhu, Menglong and Zhmoginov, Andrey and Chen, Liang-Chieh},
  booktitle={Proceedings of the IEEE Conference on Computer Vision and Pattern Recognition (CVPR)},
  pages={4510--4520},
  year={2018}
}